\chapter{測度分解定理}
\label{ch:disintegration}

接着補題（補題~\ref{lem:wp-gluing}）で必要なのは，
二変数の分布 $\pi$ を，共通の変数 $y$ ごとの分布 $\pi_y$ に分ける次の形である．
北別府悠『最適輸送理論入門〜兼非標準的関数解析入門〜』Theorem 5.3.1 の
証明（提供資料）を参考に，compact な場合の構成から一般の Polish 空間へ拡張する．
証明には第\ref{ch:radon-nikodym}章の Radon--Nikodym の定理と，
第\ref{ch:riesz-representation}章の Riesz の表現定理を用いる．

\section{主張}

\begin{theorem}{測度分解定理（接着補題で用いる形）}{meas-disintegration}
  $\mathcal{X},\mathcal{Y}$ を Polish 空間，$\pi\in\Prob(\mathcal{X}\times\mathcal{Y})$ とし，
  第2周辺分布 $\nu\in\Prob(\mathcal{Y})$ を
  \[
    \nu(B):=\pi(\mathcal{X}\times B)\qquad(B\in\Borel(\mathcal{Y}))
  \]
  で定める．
  このとき，確率測度の族 $(\pi_y)_{y\in\mathcal{Y}}$，$\pi_y\in\Prob(\mathcal{X})$ で，
  次をみたすものが存在する．
  \begin{enumerate}
    \item 各 Borel 集合 $A\in\Borel(\mathcal{X})$ に対して $y\mapsto\pi_y(A)$ は Borel 可測である．
    \item 任意の Borel 集合 $A\in\Borel(\mathcal{X})$，$B\in\Borel(\mathcal{Y})$ に対して
      \[
        \pi(A\times B)=\int_B\pi_y(A)\,\d\nu(y).
      \]
  \end{enumerate}
  さらに，$(\pi'_y)_{y\in\mathcal{Y}}$ も条件 1，2 をみたすならば，
  ある $\nu$-零集合 $N$ が存在して，$y\notin N$ ならば $\pi_y=\pi'_y$ である．
\end{theorem}

条件 1，2 をみたす族 $(\pi_y)$ を $\pi$ の $\nu$ に関する\textbf{分解}という．
証明は，第1の空間が compact な場合の存在，一般の場合の存在，一意性の順に与える．
以下，固定した $x_0\in\mathcal{X}$ に対する Dirac 測度 $\delta_{x_0}$ を，
例外集合上の分解として用いる．

\section{第1の空間が compact な場合}

一般の場合への拡張で用いるため，ここでは $\pi$ を確率測度に限らず有限測度とする．

\begin{lemma}{compact な場合の分解}{dis-compact}
  $K$ を空でない compact 距離空間，$\mathcal{Y}$ を距離空間とし，
  $\pi$ を $K\times\mathcal{Y}$ 上の零でない有限 Borel 測度，
  $\nu(B):=\pi(K\times B)$ とする．
  このとき，$\pi_y\in\Prob(K)$ の族で，各 $A\in\Borel(K)$ について
  $y\mapsto\pi_y(A)$ が Borel 可測であり，
  \[
    \pi(A\times B)=\int_B\pi_y(A)\,\d\nu(y)
    \qquad(A\in\Borel(K),\ B\in\Borel(\mathcal{Y}))
  \]
  をみたすものが存在する．
\end{lemma}

\begin{proof}
  \textbf{1. 積分の候補を作る．}
  非負の $g\in C(K)$ に対し
  \[
    \nu_g(B):=\int_{K\times B}g(x)\,\d\pi(x,y)
    \qquad(B\in\Borel(\mathcal{Y}))
  \]
  とおく．単調収束定理より $\nu_g$ は有限測度であり，
  $\nu_g(B)\leq\lVert g\rVert_\infty\nu(B)$ より，
  $\nu(B)=0$ ならば $\nu_g(B)=0$ である．
  Radon--Nikodym の定理（定理~\ref{thm:meas-radon-nikodym}，$2\Rightarrow3$）から，
  非負で $\nu$-可積分な Borel 可測関数 $l_g$ で
  $\nu_g(B)=\int_Bl_g\,\d\nu$ となるものを取る．
  ただし $l_0:=0$，$l_1:=1$ と選ぶ（$\nu_0=0$，$\nu_1=\nu$ だから選べる）．
  一般の $f\in C(K)$ には $l_f:=l_{f^+}-l_{f^-}$ と定める
  （$f^\pm$ は注意~\ref{rem:meas-compact-continuous}）．
  $f\geq0$ なら $f^-=0$ なので，$l_f=l_{f^+}\geq0$ が至るところで成り立つ．
  また積分の線形性より，すべての $f\in C(K)$ と $B$ で
  \[
    \int_Bl_f\,\d\nu=\int_{K\times B}f(x)\,\d\pi(x,y)
  \]
  である．

  $1$ を含む可算な稠密 $\mathbb{Q}$-ベクトル部分空間 $\mathcal{C}\subset C(K)$ を取る
  （定理~\ref{thm:meas-ck-separable}）．
  $f,g\in\mathcal{C}$，$a,b\in\mathbb{Q}$ について，$l_{af+bg}$ と $al_f+bl_g$ は
  上の積分公式により任意の $B$ で同じ積分値をもつので，
  密度の一意性（補題~\ref{lem:rn-density-uniqueness}）より $\nu$-a.e. に一致する．
  この条件は可算個なので，各条件の例外集合の合併として，
  Borel 集合である $\nu$-零集合 $N$ で，$y\notin N$ ならば
  \[
    l_{af+bg}(y)=a\,l_f(y)+b\,l_g(y)\qquad(f,g\in\mathcal{C},\ a,b\in\mathbb{Q})
  \]
  となるものが取れる．
  $y\notin N$，$f\in\mathcal{C}$ とし，有理数 $r>\lVert f\rVert_\infty$ を取ると，
  $r\pm f\geq0$ かつ $r\pm f\in\mathcal{C}$ より $r\pm l_f(y)=l_{r\pm f}(y)\geq0$ である．
  $r$ を $\lVert f\rVert_\infty$ に近づけて $\abs{l_f(y)}\leq\lVert f\rVert_\infty$ を得る．

  \textbf{2. 積分の候補から測度を得る．}
  $y\notin N$ とする．$f\mapsto l_f(y)$ は $\mathcal{C}$ 上で $\mathbb{Q}$-線形であり，
  1 より $\abs{l_f(y)-l_g(y)}=\abs{l_{f-g}(y)}\leq\lVert f-g\rVert_\infty$ をみたす．
  $f\in C(K)$ に対し，$\lVert f_k-f\rVert_\infty\to0$ となる $f_k\in\mathcal{C}$ を取ると
  $(l_{f_k}(y))_k$ は Cauchy 列なので，
  \[
    L_y(f):=\lim_{k\to\infty}l_{f_k}(y)
  \]
  と定める．極限は近似列の選び方によらず，$f\in\mathcal{C}$ なら $L_y(f)=l_f(y)$ である．
  加法性は極限で保たれ，有理係数を実数に近似すれば実線形性も従う．
  また $L_y(1)=1$，$\abs{L_y(f)}\leq\lVert f\rVert_\infty$ である．
  $f\geq0$ なら，$m:=\lVert f\rVert_\infty$ とおくと $0\leq m-f\leq m$ より
  $\lVert m-f\rVert_\infty\leq m$ なので，
  \[
    L_y(f)=m-L_y(m-f)\geq0.
  \]
  したがって Riesz の表現定理（定理~\ref{thm:meas-riesz-representation}）により
  $L_y(f)=\int_Kf\,\d\pi_y$（$f\in C(K)$）となる $\pi_y\in\Prob(K)$ が得られる．
  $y\in N$ では $\pi_y:=\delta_{x_0}$ とする（$x_0\in K$）．

  \textbf{3. 連続関数に対する積分公式．}
  $f\in C(K)$ と上の近似列 $f_k\in\mathcal{C}$ を取る．
  $y\notin N$ で $\int_Kf\,\d\pi_y=\lim_kl_{f_k}(y)$ であり，$N$ は Borel 集合なので，
  $y\mapsto\int_Kf\,\d\pi_y$ は Borel 可測である．
  さらに $\abs{l_{f_k}}\leq\lVert f_k\rVert_\infty\leq\lVert f\rVert_\infty+1$（$k$ が大きいとき，$N$ の外）
  なので，有界収束定理を両辺に用いて
  \[
    \int_B\left(\int_Kf\,\d\pi_y\right)\d\nu(y)
    =\lim_k\int_Bl_{f_k}\,\d\nu
    =\lim_k\int_{K\times B}f_k(x)\,\d\pi(x,y)
    =\int_{K\times B}f(x)\,\d\pi(x,y).
  \]

  \textbf{4. Borel 集合への拡張．}
  空でない閉集合 $A\subset K$ に対し，補題~\ref{lem:meas-closed-approx}の
  $\phi_k\downarrow\mathbf{1}_A$ を 3 の $f$ として有界収束定理を用いると，
  $y\mapsto\pi_y(A)$ は可測で，主張の積分公式が成り立つ（$A=\emptyset$ では明らか）．
  可測性と積分公式をみたす $A\in\Borel(K)$ の全体は $\lambda$-系である：
  $K$ を含み（$\pi_y(K)=1$，$\pi(K\times B)=\nu(B)$），補集合については
  $\pi_y(K\setminus A)=1-\pi_y(A)$ と $\pi((K\setminus A)\times B)=\nu(B)-\pi(A\times B)$ を，
  互いに素な可算合併については単調収束定理を用いればよい．
  閉集合の全体は $\Borel(K)$ を生成する $\pi$-系なので，
  Dynkin の定理（定理~\ref{thm:meas-pi-lambda}）によりすべての $A\in\Borel(K)$ で成り立つ．
\end{proof}

\section{一般の Polish 空間の場合}

\begin{proof}[定理~\ref{thm:meas-disintegration}の存在の証明]
  \textbf{1. compact 集合による分割．}
  第1周辺分布 $\mu(A):=\pi(A\times\mathcal{Y})$ は Prokhorov の定理
  （定理~\ref{thm:meas-prokhorov}，$1\Rightarrow2$）より tight である
  （$\{\mu\}$ 内の列は定数列で $\mu$ に弱収束する）．
  $K_0:=\emptyset$ とし，tightness で得た compact 集合と $K_{n-1}$ の合併を取って，
  $K_{n-1}\subset K_n$，$\mu(\mathcal{X}\setminus K_n)\leq1/n$ となる compact 集合の列を作る．
  $K:=\bigcup_nK_n$ とおけば $\mu(\mathcal{X}\setminus K)=0$ なので，
  任意の $A\in\Borel(\mathcal{X})$，$B\in\Borel(\mathcal{Y})$ について
  \[
    \pi(A\times B)=\pi((A\cap K)\times B).
  \]
  $E_n:=K_n\setminus K_{n-1}$，$\nu_n(B):=\pi(E_n\times B)$ とおけば，
  $E_n$ は $K$ を互いに素に分割するので，可算加法性より $\sum_n\nu_n(B)=\nu(B)$ である．

  \textbf{2. 各部分での分解．}
  $\nu_n\neq0$ の各 $n$ では，$\pi$ の $E_n\times\mathcal{Y}$ への制限
  $C\mapsto\pi(C\cap(E_n\times\mathcal{Y}))$ を $K_n\times\mathcal{Y}$ 上の有限測度とみて
  補題~\ref{lem:dis-compact}を適用する
  （閉部分空間の Borel 集合は周囲の Borel 集合である．注意~\ref{rem:meas-subspace-borel}）．
  得た $K_n$ 上の確率測度を，$A\in\Borel(\mathcal{X})$ に $A\cap K_n$ の測度を割り当てて
  $\mathcal{X}$ 上に延長し，$\rho_y^{(n)}$ と書く．このとき
  $y\mapsto\rho^{(n)}_y(A)$ は Borel 可測で，
  \[
    \pi((A\cap E_n)\times B)=\int_B\rho_y^{(n)}(A)\,\d\nu_n(y)
    \qquad(A\in\Borel(\mathcal{X}),\ B\in\Borel(\mathcal{Y}))
  \]
  である．$\nu_n=0$ の $n$ では $\rho_y^{(n)}:=\delta_{x_0}$ とすれば，上式は両辺 $0$ で成り立つ．

  \textbf{3. 重みを付けて足し合わせる．}
  $\nu_n(B)\leq\nu(B)$ なので，Radon--Nikodym の定理（定理~\ref{thm:meas-radon-nikodym}）から
  非負の Borel 可測な密度 $h_n:=\d\nu_n/\d\nu$ を取る．
  単調収束定理より $\int_B\sum_nh_n\,\d\nu=\sum_n\nu_n(B)=\nu(B)=\int_B1\,\d\nu$ なので，
  密度の一意性（補題~\ref{lem:rn-density-uniqueness}）より $\sum_nh_n=1$ が $\nu$-a.e. に成り立つ．
  その例外集合 $N_0:=\{\sum_nh_n\neq1\}$ は Borel 集合である．そこで
  \[
    \pi_y(A):=
    \begin{cases}
      \sum_nh_n(y)\,\rho_y^{(n)}(A) & (y\notin N_0),\\
      \delta_{x_0}(A) & (y\in N_0)
    \end{cases}
  \]
  と定める．$y\notin N_0$ では，非負係数の和が $1$ の確率測度の可算和なので，
  単調収束定理（数列の場合）より $\pi_y$ は確率測度である．
  各項の可測性から条件 1 が成り立つ．
  系~\ref{cor:rn-density-integral}と単調収束定理より
  \[
    \begin{aligned}
      \int_B\pi_y(A)\,\d\nu(y)
      &=\sum_n\int_B\rho_y^{(n)}(A)\,h_n(y)\,\d\nu(y)
      =\sum_n\int_B\rho_y^{(n)}(A)\,\d\nu_n(y)\\
      &=\sum_n\pi((A\cap E_n)\times B)=\pi((A\cap K)\times B)=\pi(A\times B).
    \end{aligned}
  \]
  よって条件 2 も成り立つ．
\end{proof}

\section{一意性}

\begin{proof}[定理~\ref{thm:meas-disintegration}の一意性の証明]
  $(\pi_y)$，$(\pi'_y)$ がともに条件 1，2 をみたすとする．
  $A\in\Borel(\mathcal{X})$ を固定すると，$y\mapsto\pi_y(A)$，$y\mapsto\pi'_y(A)$ は
  $[0,1]$ に値をとる可測関数であり，条件 2 よりすべての $B$ で同じ積分値をもつ．
  密度の一意性（補題~\ref{lem:rn-density-uniqueness}）より，
  ある $\nu$-零集合 $N_A$ の外で $\pi_y(A)=\pi'_y(A)$ である．

  補題~\ref{lem:meas-countable-generator}により，
  $\Borel(\mathcal{X})$ を生成する可算な $\pi$-系 $\mathcal{E}$ を取り，
  $N:=\bigcup_{A\in\mathcal{E}}N_A$ とおく．$N$ は $\nu$-零集合である．
  $y\notin N$ では $\pi_y,\pi'_y$ は $\mathcal{E}$ 上で一致する確率測度なので，
  一致定理（定理~\ref{thm:meas-uniqueness}）より $\pi_y=\pi'_y$ である．
\end{proof}
